% !TEX program = pdflatex
% Fordítás:  latexmk -pdf main.tex   (a biber-t a latexmk magától lefuttatja)
\documentclass[a4paper,12pt,oneside]{article}

% ---------------------------------------------------------------------
% Nyelv és betűkészlet
% ---------------------------------------------------------------------
\usepackage[T1]{fontenc}
\usepackage[utf8]{inputenc}
\usepackage[magyar]{babel}
\usepackage{amsmath}
\usepackage{mathptmx}   % Times stílusú matematikai betűk
\usepackage{tgtermes}   % TeX Gyre Termes: Times New Roman-klón, ő/ű is
\usepackage{microtype}

% ---------------------------------------------------------------------
% Formai követelmények (TDK 2.2)
% ---------------------------------------------------------------------
% A4, margók: bal 3 cm, jobb 2 cm, felső 3 cm, alsó 3 cm
\usepackage[a4paper,left=3cm,right=2cm,top=3cm,bottom=3cm]{geometry}

% 1,5-es sorköz
\usepackage{setspace}
\onehalfspacing

% Fejezetcímek: 16 pt (fejezet), 14 pt (alfejezet), 12 pt (alalfejezet)
\usepackage{titlesec}
\titleformat{\section}
  {\normalfont\fontsize{16}{20}\selectfont\bfseries}{\thesection.}{0.6em}{}
\titleformat{\subsection}
  {\normalfont\fontsize{14}{18}\selectfont\bfseries}{\thesubsection.}{0.6em}{}
\titleformat{\subsubsection}
  {\normalfont\normalsize\bfseries}{\thesubsubsection.}{0.6em}{}
\titlespacing*{\section}{0pt}{3.0ex plus 1ex minus .2ex}{1.5ex plus .2ex}
\titlespacing*{\subsection}{0pt}{2.5ex plus 1ex minus .2ex}{1.0ex plus .2ex}

% ---------------------------------------------------------------------
% Ábrák, táblázatok, hivatkozások
% ---------------------------------------------------------------------
\usepackage{graphicx}
\usepackage{booktabs}
\usepackage{csquotes}
\usepackage[backend=biber,style=numeric-comp,sorting=none]{biblatex}
\addbibresource{refs.bib}

\usepackage[hidelinks]{hyperref}   % nyomtatásban nem látszó linkek

% ---------------------------------------------------------------------
% A dolgozat címe (a címlapon és a PDF adataiban is ez szerepel)
% ---------------------------------------------------------------------
\newcommand{\dolgozatcim}{Autonóm járművezetés megerősítéses tanulással\\a CARLA szimulátorban}
\newcommand{\dolgozatalcim}{Kamera- és LiDAR-autoenkóderre épülő\\Soft Actor-Critic ágens}

\hypersetup{
  pdftitle   = {Autonóm járművezetés megerősítéses tanulással a CARLA szimulátorban},
  pdfauthor  = {Angi Dávid},
  pdfsubject = {TDK dolgozat},
}

\begin{document}

% =====================================================================
% Címlap
% =====================================================================
\begin{titlepage}
  \singlespacing
  \centering

  % Egyetem, kar, tanszék -- a lap tetején, középen
  {\fontsize{14}{18}\selectfont\bfseries Debreceni Egyetem\par}
  \vspace{4pt}
  {\large Informatikai Kar\par}
  {\large Adattudomány és Vizualizáció Tanszék\par}
  \vspace{10pt}
  \rule{\textwidth}{0.6pt}

  % A pályamunka címe -- a felső harmadban, középen
  \vspace{0.06\textheight}
  {\fontsize{16}{22}\selectfont\bfseries \dolgozatcim\par}
  \vspace{14pt}
  {\fontsize{14}{18}\selectfont \dolgozatalcim\par}
  \vspace{22pt}
  {\normalsize\scshape Tudományos Diákköri Dolgozat\par}

  \vspace{\stretch{3}}

  % Témavezető (bal) és szerző (jobb) -- az alsó harmadban
  \begin{minipage}[t]{0.48\textwidth}
    \raggedright
    {\bfseries Témavezető:}\\[6pt]
    Mészáros László\\
    egyetemi adjunktus % TODO: pontos beosztás
  \end{minipage}%
  \hfill
  \begin{minipage}[t]{0.48\textwidth}
    \raggedleft
    {\bfseries Készítette:}\\[6pt]
    Angi Dávid\\
    programtervező informatikus BSc % TODO: szak, évfolyam
  \end{minipage}

  \vspace{\stretch{1}}

  % Benyújtás helye és éve -- a lap alján, középen
  \rule{\textwidth}{0.6pt}\\[8pt]
  {\large Debrecen, 2026\par}
\end{titlepage}
\setcounter{page}{2}

% =====================================================================
% Tartalomjegyzék
% =====================================================================
\tableofcontents
\newpage

% =====================================================================
% Törzsszöveg
% =====================================================================
\section{Bevezetés}
\label{sec:bevezetes}

Az autonóm járművezetés egyik ígéretes megközelítése a megerősítéses
tanulás (reinforcement learning, RL)~\cite{sutton2018rl}, amelyben az ágens
a környezettel való interakció során, jutalomjelek alapján tanulja meg a
vezetési stratégiát~\cite{kiran2022survey,kendall2019learning}. A valós
járművön történő tanítás költséges és veszélyes, ezért a kísérleteket a
nyílt forráskódú CARLA szimulátorban~\cite{dosovitskiy2017carla} végezzük.

% TODO: motiváció, célkitűzés, a dolgozat felépítése

\section{Kapcsolódó munkák}
\label{sec:kapcsolodo}

% TODO: RL alapú vezetés, látens reprezentációk, a kiinduló CARLA-SB3 környezet
Munkánk kiindulópontja a CARLA-SB3 tanítási környezet~\cite{mate2023carlasb3},
amely egy variációs autoenkóder látens terével tanít SAC ágenst. Hasonló,
látens térben tanuló megközelítést mutat be \textcite{chen2022interpretable}.

\section{Elméleti háttér}
\label{sec:elmelet}

\subsection{Megerősítéses tanulás}
% TODO: MDP, politika, értékfüggvény

\subsection{Soft Actor-Critic}
% TODO: maximum entrópiás RL, SAC~\cite{haarnoja2018sac}
Az ágenst a Soft Actor-Critic algoritmussal~\cite{haarnoja2018sac} tanítjuk,
a Stable-Baselines3 könyvtár~\cite{raffin2021sb3} implementációját használva.

\subsection{Autoenkóderek}
% TODO: dimenziócsökkentés autoenkóderrel~\cite{hinton2006reducing}
A nagy dimenziós szenzoradatokat autoenkóderrel~\cite{hinton2006reducing}
tömörítjük kis dimenziós látens vektorrá.

\section{A javasolt rendszer}
\label{sec:rendszer}

\subsection{Szimulációs környezet}
% TODO: CARLA, Town04, útvonaltervezés (GlobalRoutePlanner)

\subsection{Érzékelés és látens reprezentáció}
% TODO: kamera-autoenkóder, LiDAR range-image gráf-autoenkóder

\subsection{Megfigyelési és akciótér}
% TODO: látens vektorok, irányszög, manőver; kormány és gáz

\subsection{Jutalomfüggvény}
% TODO: sebesség, célba érés, ütközés elkerülése

\subsection{Tanítási beállítások}
% TODO: SAC hiperparaméterek (táblázat)

\section{Kísérletek és eredmények}
\label{sec:eredmenyek}

\subsection{Kísérleti beállítások}
% TODO

\subsection{Eredmények}
% TODO: tanulási görbék, összehasonlítás a kiinduló környezettel

\subsection{Ablációs vizsgálat}
% TODO

\section{Összegzés és további munka}
\label{sec:osszegzes}

% TODO: eredmények összefoglalása, forgalom (NPC járművek) hozzáadása

% =====================================================================
% Irodalomjegyzék
% =====================================================================
\newpage
\printbibliography[heading=bibintoc,title={Irodalomjegyzék}]

\end{document}
